\appendix
\chapter{Recomandări Diagrame și Prompt-uri pentru Vizualizări PhD-Level}
\label{cap:diagrame_recomandari}

Acest capitol furnizează recomandări detaliate și prompt-uri "PhD-level" pentru crearea/îmbunătățirea diagramelor din teză. Scopul este de a standardi calitatea vizualizărilor și de a îndruma procesul de generare a imaginilor profesionale.

\section{1. Diagramă: Arhitectura Sistemului End-to-End}

\textbf{Locație recomandată}: Capitolul 5 (Implementare), după "Schema Generală a Aplicației"

\textbf{Tip}: Arhitectural diagram (block/flowchart)

\textbf{Prompt detaliat pentru generare/descriere}:

Creează o diagramă arquitecturală completă care arată fluxul datelor în StyleGenAI de la utilizator la recomandare finală. Include:

- Stânga: Utilizatorul și dispozitivul mobil (smartphone cu React Native app)
- Centru: Cloud/Backend cu trei trepte:
  1. API Gateway (Flask router) care primește requesturi HTTP POST
  2. Service Layer cu trei componente paralele:
     - Image Processing Pipeline (rembg, color extraction)
     - Recommendation Engine (outfit generation + scoring)
     - Trend Analysis Service (trend color lookup)
  3. Persistence Layer (SQLite database cu 4 tabele: Users, Clothes, Outfits, TrendColors)
- Dreapta: Săgeți de response către frontend cu JSON payload (outfit recommendation + score)
- Legenda: Colorează API Gateway cu verde, Services cu albastru, Database cu portocaliu
- Adnotări pe fiecare arrow cu tipuri de date (JWT token, FormData, JSON)
- Inclusivă: Security layer labeled "JWT Validation" între Gateway și Services

\textbf{Instrumente sugerate}: Lucidchart, Diagram.io, TikZ (LaTeX native), Figma

\textbf{Aspect critic}: Trebuie să fie evident că JWT validation este checkpoint central; Services sunt independent escalabile.

---

\section{2. Diagramă: Pipelineul de Procesare Imagine}

\textbf{Locație recomandată}: Capitolul 5, Subsecția "Formalizare Matematică"

\textbf{Tip}: Process flow diagram

\textbf{Prompt detaliat}:

Creează o diagramă care arată transformările imaginii de la raw input la output (detected color). Structura:

1. Input: Raw photo (JPG, PNG) cu un articol vestimentar în fond neuniform
2. Etapa 1 - Background Removal: Arată original vs. output al rembg (PNG transparent)
3. Etapa 2 - RGB Extraction: Arată pixelii izolaților și histograma de frecvență RGB
4. Etapa 3 - Color Quantization: Arată tranziția de la RGB continuu la paleta CSS3 discretă (show ~140 culori web standard)
5. Etapa 4 - Euclidean Distance Calculation: Formule pe diagramă pentru d(c, p_i)
6. Output: Culoare dominantă identificată și labeled cu hex code (e.g., #FF5733)

Include în lateral micile diagrame 3D ale spațiului RGB cu distanțe highlight-te din culorile articolelor.

---

\section{3. Diagramă: State Machine - Generare Outfit}

\textbf{Locație recomandată}: Capitolul 5, subsecția "Algoritm de Generare Ținutelor"

\textbf{Tip}: State diagram / flowchart

\textbf{Prompt detaliat}:

Creează o diagramă care modelează algoritmul de generare a outfit-ului ca o mașină de stare. Secvența:

1. START STATE: "Primire Request"
2. Stare 2: "Load User Clothes din DB" (Query filtered by user_id)
3. Stare 3: "Categorizare (Top/Bottom/Shoe)" - paralel proceseaza
4. Stare 4: "Generate All Combinations (nested loops)" - show O(n^3) complexity
5. Stare 5: "Calculare Scores" - paralel evaluează S(c_i, c_j) pentru fiecare combo
6. Stare 6: "Aplicare Trend Boost" - conditional: dacă culoare in trendColors, add boost
7. Stare 7: "Rank & Select Top-K" - sorteaza dupa score descending
8. OUTPUT STATE: "Return Outfit + Score" (JSON response)

Inclusiv: Decision nodes pentru validări (sunt clothes suficiente? are user favorite filters?); error states pentru edge cases (zero clothes, categorii necomplet).

---

\section{4. Diagramă: Baza de Date - ERD Detaliată (Entity-Relationship Diagram)}

\textbf{Locație recomandată}: Capitolul 5, subsecția "Securitate și Baza de Date"

\textbf{Tip}: ER diagram

\textbf{Prompt detaliat}:

Creează un ERD profesional care arată:

Tabele (fiecare cu subset de coloane):
- USERS (user_id PK, email UNIQUE, password_hash, created_at, updated_at)
- CLOTHES (cloth_id PK, user_id FK->USERS, category ENUM, dominant_color VARCHAR, image_path, style, created_at)
- OUTFITS (outfit_id PK, user_id FK->USERS, top_id FK->CLOTHES, bottom_id FK->CLOTHES, shoe_id FK->CLOTHES, compatibility_score DECIMAL, liked BOOLEAN, created_at)
- TREND_COLORS (trend_id PK, rgb_value VARCHAR, season VARCHAR, year INT, popularity_score DECIMAL)

Relații:
- USERS --[1:N]--> CLOTHES (owneds)
- USERS --[1:N]--> OUTFITS (generates)
- CLOTHES --[1:N]--> OUTFITS (includes) - 3x relații (top, bottom, shoe)

Indexuri: Highlight indexuri composite pe (user_id, category) pentru CLOTHES; (user_id, created_at) pentru OUTFITS.

Constraints: NOT NULL pe campos cheie; CHECK constraints pentru enums; FOREIGN KEY cascading deletes on user delete.

---

\section{5. Diagramă: Model Conceptual de Compatibilitate Cromatică}

\textbf{Locație recomandată}: Capitolul 5, subsecția "Formalizare Matematică"

\textbf{Tip}: Conceptual/Mathematical visualization

\textbf{Prompt detaliat}:

Creează o vizualizare care explică funcția de compatibilitate $S(c_1, c_2) = \alpha H + \beta M + \gamma N + \delta T$:

Arată 4 paneluri, fiecare pentru o componentă:

1. Panel H (Harmonia): Arată color wheel HSV cu două culori la 180° (complementare, H=1.0) și la alte unghiuri (H decreasing). Include formula: $H(c_1, c_2) = 1 - |(\text{hue}(c_1) - \text{hue}(c_2)) - 180°| / 180°$

2. Panel M (Monochromatic): Arată spectru RGB cu două culori apropiate (M=1.0) și departe (M→0). Formula: $M(c_1, c_2) = 1 - d(c_1, c_2) / d_{max}$

3. Panel N (Neutral): Arată legenda neutru (alb, negru, gri) vs. accent culori. Formula: $N(c_1, c_2) = \text{bonus dacă oricare din culori este neutru}$

4. Panel T (Trend): Arată calendar cu sezoane și culori populare. Formula: $T(c_1, c_2) = \text{boost dacă ambele culori sunt in trend colors}$

Bottom: Exemplu concret de outfit (3 imagini) cu breakdown de scor: "Alb (neutral) + Bleumarin (armonie 85°, monocromatism 70%) + Negru (neutral) = Total 0.87"

---

\section{6. Diagramă: Fluxul UI Mobile - Wireflow cu Data Binding}

\textbf{Locație recomandată}: Capitolul 5, subsecția "Dezvoltarea Interfeței Mobile"

\textbf{Tip}: Wireflow diagram

\textbf{Prompt detaliat}:

Creează un wireflow care arată navigația în aplicație și binding-ul datelor:

Ecrane:
1. Login Screen → POST /auth/login → receive JWT token → store în AsyncStorage
2. Wardrobe Screen (Tab 1) → GET /clothes?user_id={userId} → display grid de imagini
   - Button "Add Photo" → modal camera/gallery → POST /clothes/upload → success toast
3. Generator Screen (Tab 2) → POST /get_suggestion (cu FormData imagini selectate) → receive Outfit JSON → render OutfitCard
   - Like button → POST /outfits/{id}/like → update UI (heart icon becomes red)
   - Swipe left → DELETE /outfits/{id}
4. History Screen (Tab 3) → GET /outfits/history → display cards din array response

Fiecare ecran arată:
- Input/output state (ce data intră, ce iese)
- Headers HTTP necesari (Authorization: Bearer {token})
- Response schemas în JSON snippet-uri

---

\section{7. Diagramă: Comparație Performanță - Timeline Chart}

\textbf{Locație recomandată}: Capitolul 6 (Testare), subsecția "Detalii Măsurători"

\textbf{Tip}: Timeline/bar chart

\textbf{Prompt detaliat}:

Creează o visualizare care compară timpii de răspuns ai endpoint-urilor:

1. Axa X: Endpoint-uri (/health, /get_suggestion, /web_outfit, /outfits/history, /outfits/{id}/like)
2. Axa Y: Timp (ms), scară 0-1000
3. Pentru fiecare endpoint, arată:
   - Min, Max, Mean ca bar chart cu eroare bars
   - Scatter plot overlay showing individual measurements

Culori:
- Verde: sub 50 ms (fast)
- Galben: 50-200 ms (acceptable)
- Portocaliu: 200-500 ms (slower but acceptable)
- Roșu: >500 ms (slow, needs optimization)

Linie orizontală de referință: 1000 ms (user frustration threshold).

---

\section{8. Diagramă: Rezultate Testare - Heatmap de Acuratețe}

\textbf{Locație recomandată}: Capitolul 6, subsecția "Evaluarea Calității"

\textbf{Tip}: Heatmap

\textbf{Prompt detaliat}:

Creează o heatmap care arată acuratețea clasificării pe categorii vestimentare:

Matrice 3x3 (Predicted vs. Actual):
- Rânduri: Actual (Top, Bottom, Shoe)
- Coloane: Predicted (Top, Bottom, Shoe)
- Celule: Count + percentage

Colorează celule după intensitate:
- Verde închis: 95-100% (perfect)
- Verde deschis: 80-94% (good)
- Galben: 60-79% (fair)
- Roșu: <60% (poor)

Inclusiv: Precision, Recall, F1-score în margini.

---

\section{9. Diagramă: Corelație Score Algoritm vs. Evaluare Umană}

\textbf{Locație recomandată}: Capitolul 6, subsecția "Evaluarea Calității"

\textbf{Tip}: Scatter plot cu trend line

\textbf{Prompt detaliat}:

Creează un scatter plot care arată corelația între scorul algoritmic și notele umane:

- Axa X: Scor Algoritm (0.6-1.0)
- Axa Y: Notă Umană Medie (1-5)
- Fiecare punct: un outfit testat (25 puncte total)
- Trend line: Linear regression cu r=0.78 labeled
- Confidence band (±95%) în jur de linie

Include: Textbox cu statistici (r-value, p-value, equation).

---

\section{10. Diagramă: SWOT Analysis - Canvas Visual}

\textbf{Locație recomandată}: Capitolul 3 (Studiu Bibliografic), după tabel SWOT

\textbf{Tip}: SWOT quadrant diagram

\textbf{Prompt detaliat}:

Creează o SWOT analysis visualization (2x2 quadrant):

- Top-Left (Strengths): Transparency, Speed, Usability, Modularity - highlight cu verde
- Top-Right (Weaknesses): MVP stage, Scaling challenges, Hardcoded trends, Basic ML - highlight cu roșu
- Bottom-Left (Opportunities): E-Commerce, AR, RL personalization, Global expand - highlight cu albastru
- Bottom-Right (Threats): Competition (Acloset, Whering), Saturation, Social fashion shifts, BigTech - highlight cu portocaliu

Fiecare quadrant: bullet points cu iconuri reprezentative.

---

\section{11. Diagramă: Regresia Liniară - Performance vs. Image Size}

\textbf{Locație recomandată}: Capitolul 6 (Testare), investigații suplimentare

\textbf{Tip}: Scatter plot cu regression curve

\textbf{Prompt detaliat}:

Creează un plot care arată cum variază performance în funcție de dimensiunea imaginii încărcate:

- Axa X: Image size (KB) 100-5000
- Axa Y: Response time (ms) 300-900
- Puncte: Măsurători reale din 30 teste
- Linie: Linear regression cu ecuația afișată
- Inclusiv: Outliers identify-ați și analizați

Mesaj: Are o corelaț strongly positive, sugerând că optimization de image compression ar ajuta.

---

\section{12. Diagramă: Roadmap Viitor - Timeline cu Milestones}

\textbf{Locație recomandată}: Capitolul 8 (Concluzii), subsecția "Direcții Viitoare"

\textbf{Tip}: Gantt chart / timeline

\textbf{Prompt detaliat}:

Creează un roadmap visual cu faze viitoare de dezvoltare:

Timeline: 6-24 luni

Faze:
1. Q2 2026: Beta launch, public UAT (4 săptămâni)
2. Q3 2026: Feedback integration, bug fixes, security audit (8 săptămâni)
3. Q4 2026: AR Try-On MVP integration (12 săptămâni, paralel cu faza 2)
4. Q1 2027: E-Commerce partnerships, API integrations (8 săptămâni)
5. Q2 2027: RL personalization module (16 săptămâni, complexă)
6. Q3 2027: Launch v2.0 production-ready

Fiecare fază: color-coded (verde=planning, albastru=development, portocaliu=testing, roșu=launch)

---

\section{Instrumente și Alte Recomandări}

\subsection{Software/Tools}

\begin{itemize}
    \item \textbf{Architecture/Flow Diagrams}: Lucidchart, Diagram.io, Draw.io, Figma, OmniGraffle
    \item \textbf{Data Visualization}: Matplotlib (Python), Plotly, R (ggplot2), Tableau
    \item \textbf{LaTeX Native}: TikZ/pgfplots (native LaTeX, consistency garantată)
    \item \textbf{Screenshots/UI}: Figma, Sketch, AdobeXD (din React Native simulator)
\end{itemize}

\subsection{Standarde de Calitate}

\begin{itemize}
    \item Rezoluție: Minimum 300 DPI pentru print
    \item Format: PDF preferred pentru vector graphics (scalable)
    \item Culori: Palette profesională, accessible (colorblind-friendly)
    \item Font: Consistent cu teză (Cambria sau Times New Roman pentru math)
    \item Legenda: Obligatorie; fontsize readable
    \item Caption: Descriptiv, cu referință la text
\end{itemize}

\subsection{Checklist Pre-Submission}

Pentru fiecare diagramă, verifică:
\begin{enumerate}
    \item Data accuracy - sunt numerele/statisticile corecte și din teză?
    \item Readability - se vede clar pe format A4 în print?
    \item Alignment - sunt tous diagrama și text align-ate bine în pagină?
    \item Labeling - sunt todas componentele labeled și explicite?
    \item Citation - este diagramă și caption-ul citat în text capitolul respectiv?
\end{enumerate}

